\section{Discussion}\label{sec:discussion}
\paragraph{What the evidence supports} On the problems solved by third-party code, cross-validated selection from a small pool of classical models found the oracle or a near-oracle member (H1a; a fixed choice of one full-field Gaussian process does almost as well), was not worse than always using a Gaussian process (H1b; nearly guaranteed because that model is in the pool), raised the basis flag only in cases where a full-field model was better but with low recall (H3), and gave usable forecasts of the error at four times the budget when starting from 64 runs on the first eight problems (H2). The tool was then frozen and run on eight further problems under criteria and predictions committed beforehand (\cref{sec:confirm}): selection quality, the comparison with always using one Gaussian process, the classification of whether a target is reached and the comparison with the tuned network held again, which supports these four claims beyond the eight problems on which the tool was developed. Against a tuned network the tool met the criterion of H4 only at its limit, and that network is probably under-trained. On real measured data with scalar, noisy outputs, all four descriptive criteria were met, although two of the five R2 successes are ties and the scalar outputs make the POD rank one, so these data do not test the central design of the tool.
\paragraph{What it does not support} The benchmark finding that the reduced-basis group wins at large budgets did \emph{not} hold on \TpNotClear{} of the \TpN{} problems (H5 \TpVerdict), so a fixed recommendation ``use a reduced-basis Gaussian process'' is not safe; the case for the tool is that it finds out from the data. The tool did not match the benchmark's own selection among five methods that include networks (H1c \SthreeVerdict): where networks win, the default pool, which has none, loses. The forecast was too optimistic from 256 runs on the first eight problems and from both budgets on the further eight (within a factor two in only \SvpBoneShare{} of \SvpBoneN{} cases over all sixteen problems), and its range is not an uncertainty interval; we do not claim that the forecast is validated. The basis flag raised no alarm on the further problems, so its usefulness rests on one batch of flags. We did not isolate marginal-likelihood fitting as the mechanism behind the Gaussian-process advantage, and the pilot diagnostic that we had hoped would predict the winner from 32 runs was not confirmed (Supplement, Sec.~\ref{sec:S-earlier}).
\paragraph{Limitations}
\begin{enumerate}
\item \textbf{Simulated problems.} All problems (the \NAllProblems{} of the study and the eight further ones) are simulated; the sixteen problems solved by third-party code come from eight families (two levels each) of one finite-element library. The real-data check uses two datasets with scalar, noisy outputs, which violate assumption A2; it says that the tool behaves sensibly on them, not that it works for noisy field data.
\item \textbf{Design knowledge.} The pool, the kernel-ridge grids and the criteria were designed with knowledge of the 36 earlier problems, so the problems solved by third-party code are out-of-sample only in part, and they were chosen by us after the earlier results were known, so ``independent'' refers to the solver, not to the choice of problems. The eight further problems were run with the tool frozen and with criteria and predictions committed first, but they too were chosen by us, after the first results were known; the smoke tests of the run scripts used a problem that had already been seen. Smoke tests of the tool on three of the eight problems preceded the commit of the criteria; they showed test errors of some pool members and exposed a numerical bug in the leave-one-out computation, which was fixed before the evaluation runs (deviation V9). The text of the criteria was not changed afterwards, but its statement that no error had been seen was wrong when written.
\item \textbf{Few independent units.} Cells within a problem are not independent, so the intervals in \cref{tab:criteria}, which treat 24 cells as independent, are optimistic. Seeds share the pool and test set and nest their training subsets.
\item \textbf{Forecast.} The forecast is a power-law extrapolation without guarantee; it was too low by more than a factor two in \BoneUnder{} of \BoneN{} cases, and the pre-specified criterion was met with a confidence interval that includes the threshold.
\item \textbf{Networks.} The default pool contains no network, the tool never recommends one, and nothing in its report warns that a network might do better. A network can be added to the pool as an option, which changed the result on two of the eight further problems at several times the cost (\cref{sec:confirm}). The tuned network used 16 trials in a restricted space (depth 2--4, width 64--256) on one core; a variant with 48 trials, depth up to five and width up to 512 was better in most problems and changed the conclusion little, but a more heavily tuned or different architecture could do better.
\item \textbf{Cost.} Fifteen of the sixteen solvers take milliseconds per run, so on these problems the tool costs more than the simulations it replaces. One problem takes \SvETsim{} seconds per run, but it is the first plate problem on a finer mesh, so it tests the cost accounting and not selection; there the surrogate was cheaper than the solver after about \SvEBreakTwofiftysix{} predictions. Simulations that take minutes or hours were not tested, and that is where the case for the tool rests.
\item \textbf{Baselines and weak test problems.} We compared with a hindsight-fixed member, a leave-one-problem-out fixed member, the benchmark's five methods and a tuned network; we did not compare with a generic stacking or AutoML selector, a multi-output Gaussian process with automatic relevance determination, or an operator-learning method (DeepONet, FNO) with a competent tuning budget, and not on a public benchmark such as PDEBench. Two of the eight further problems (undamped Helmholtz) are ones on which no model is usable, and the tool has no warning for that case.
\item \textbf{Scope of the tool.} It assumes exchangeable runs on a common output grid, deterministic outputs, and a moderate number of inputs (only \NHighP{} of our \NAllProblems{} problems have ten or more).
\item \textbf{Literature.} Our search for work on automated surrogate selection and on learning-curve extrapolation was limited and mostly at abstract level, so related work may be incomplete.
\item \textbf{Process.} Hypotheses and code were committed before the runs they concern, but the history is our own, not an independent registry; several stages were designed after seeing earlier results (labelled post hoc); the independent reviews were performed by AI agents, not human reviewers.
\end{enumerate}
\paragraph{Ethics and disclosure} No human or animal subjects were involved. The two real datasets are public (CC BY 4.0). The AI involvement is disclosed in the acknowledgments.

